\documentclass[12pt,a4paper]{amsart}

\usepackage[a4paper,textwidth=160mm,textheight=215mm]{geometry}
\usepackage{amsfonts,amssymb,amsthm,amsmath,mathrsfs}
\PassOptionsToPackage{no-math}{fontspec}
\usepackage{xeCJK}
\usepackage[hidelinks]{hyperref}

\setCJKmainfont[BoldFont={SimHei},ItalicFont={KaiTi}]{SimSun}
\setCJKsansfont{SimHei}
\setCJKmonofont{SimSun}

\linespread{1.30}
\setlength{\parindent}{2em}
\numberwithin{equation}{section}

\theoremstyle{plain}
\newtheorem{theorem}{定理}[section]
\newtheorem{proposition}[theorem]{命题}
\newtheorem{lemma}[theorem]{引理}
\newtheorem{corollary}[theorem]{推论}
\theoremstyle{definition}
\newtheorem{remark}[theorem]{注}

\renewcommand{\abstractname}{摘要}
\renewcommand{\keywordsname}{关键词}
\renewcommand{\subjclassname}{\textup{2020} 年数学主题分类}
\renewcommand{\proofname}{证明}
\renewcommand{\refname}{参考文献}

\newcommand{\HS}{\mathrm{HS}}
\newcommand{\tr}{\operatorname{tr}}
\newcommand{\dd}{\,\mathrm{d}}
\newcommand{\eps}{\varepsilon}

\begin{document}

\title[模型算子的极点谱]{条件哥德巴赫方法中\\ 模型算子的极点谱}

\author{}

\date{}

\subjclass[2020]{主分类 11M26；次分类 11M06, 47A75}

\keywords{哥德巴赫猜想，模型算子，哈代空间，欧拉乘积，迹类}

\begin{abstract}
在 Li 与 Liu 的条件哥德巴赫方法中，加权哥德巴赫误差被精确表示为一个作用在
哈代空间上的算子 $T_N=\mathcal P_N-\mathcal D_N-\mathcal K_N$ 的二次型。本文
确定模型算子 $\mathcal K_N$ 的首阶谱行为。我们证明精确分解
$\mathcal K_N=-W_N+R_N$，其中 $W_N$ 是正积分算子，其核用 \cite{LL} 中的因子
$X_1$ 以闭式写出；$R_N$ 是希尔伯特--施密特算子，满足
$\|R_N\|_{\HS}\ll(\log\log(3N))^2$。我们证明
$\|W_N\|=\bigl(\pi/e^3+o(1)\bigr)\log N$；由此 $\mathcal K_N$ 有一个负特征值
$-\pi e^{-3}\log N+o(\log N)$，且 $\|\mathcal K_N\|_{\HS}\ll(\log N)^{3/2}$。
证明不使用任何未证假设。
\end{abstract}

\maketitle

\section{引言}

\subsection*{1.1.\ 模型算子}

在 \cite{LL} 中，Li 与 Liu 证明了二元哥德巴赫猜想的一个条件形式。他们的论证
把加权哥德巴赫误差化归为竖直条带的哈代空间上一个紧自伴算子 $T_N$ 的精确二次型，
且 $T_N$ 分解为
\begin{equation}\label{eq:split}
 T_N=\mathcal P_N-\mathcal D_N-\mathcal K_N,
\end{equation}
其中 $\mathcal P_N$ 是精确的特征投影，$\mathcal D_N$ 是精确的对角算子，
$\mathcal K_N$ 是贝塔模型。本文确定 $\mathcal K_N$ 的首阶谱特征。

设 $N>e^{10}$ 为偶数。我们沿用 \cite[\S5--6]{LL} 的记号：
\[
 \alpha_N=\frac1{\log N},\qquad \beta=\frac18,\qquad
 S_N=\{s:\alpha_N<\Re s<\beta\},\qquad
 \Gamma_N=(\alpha_N+i\mathbb R)\cup(\beta+i\mathbb R),
\]
而 $\mathcal H_N=H^2(S_N)$ 是条带的哈代空间，带边界范数
\begin{equation}\label{eq:inner}
 \langle h,g\rangle_{\mathcal H_N}
 =\int_{\Gamma_N}h(s)\overline{g(s)}\,|\mathrm ds|;
\end{equation}
由 \cite[引理 5.2]{LL}，边界值在 $L^2$ 中存在，且 \eqref{eq:inner} 是
$\mathcal H_N$ 的内积。该积分是共轭对称的，我们记
$\|h\|_{\mathcal H_N}^2=\langle h,h\rangle_{\mathcal H_N}$。

基本因子为
\begin{equation}\label{eq:X1}
 X_1(s)=\frac{\Gamma\bigl(\frac{3/2-s}{2}\bigr)}
               {\Gamma\bigl(\frac{3/2+s}{2}\bigr)},
 \qquad
 P_N(w)=\prod_{\substack{p>2\\p\nmid N}}\Bigl(1+\frac{p^{-w}}{p-2}\Bigr),
 \qquad
 \mathfrak S(N)=2C_2\prod_{\substack{p\mid N\\p>2}}\frac{p-1}{p-2},
\end{equation}
其中 $C_2=\prod_{p>2}\bigl(1-(p-1)^{-2}\bigr)=0.66016\ldots$，并且
\begin{equation}\label{eq:LN}
 L_N=\log\frac N\pi,\qquad
 \mathscr A_N(w)=\mathfrak S(N)P_N(w)\zeta(1-w).
\end{equation}
模型核为
\begin{equation}\label{eq:kernel}
 \kappa_N(s,z)=\frac{X_1(s)X_1(\bar z)}{(1+s)(1+\bar z)}\,
 e^{-L_N(s+\bar z)}\mathscr A_N(s+\bar z),
 \qquad s,z\in\Gamma_N,
\end{equation}
模型算子通过下式作用于 $\mathcal H_N$：
\begin{equation}\label{eq:K}
 (\mathcal K_Nh)(s)=\int_{\Gamma_N}\kappa_N(s,z)\,h(z)\,|\mathrm dz|,
 \qquad h\in\mathcal H_N .
\end{equation}
这是 \cite[\S6.4]{LL} 中三个算子里的模型分量；它的边界因子与 \cite{LL} 相同。

对 $x\ge L_N$ 置
\begin{equation}\label{eq:Phi}
 \Phi_x(s,z)=\frac{X_1(s)X_1(\bar z)}{(1+s)(1+\bar z)}\,e^{-x(s+\bar z)},
 \qquad s,z\in\Gamma_N,
\end{equation}
并定义权重
\begin{equation}\label{eq:w}
 \varpi(x)=x-L_N+h_N,\qquad x\ge L_N,
\end{equation}
其中 $h_N$ 如下文引理~\ref{lem:value} 所定义。极点算子与余项分别是核为
\begin{equation}\label{eq:kW}
 \kappa_W(s,z)=\int_{L_N}^{\infty}\varpi(x)\,\Phi_x(s,z)\,\mathrm dx,
 \qquad
 \kappa_{R_N}(s,z)=\kappa_N(s,z)+\kappa_W(s,z)
\end{equation}
的算子；我们分别用 $W_N$ 与 $R_N$ 表示它们在 $\mathcal H_N$ 上对应的算子。

\subsection*{1.2.\ 结果陈述}

\begin{theorem}\label{thm:main}
设 $N>e^{10}$ 为偶数，$\mathcal K_N$ 为模型算子
\eqref{eq:kernel}--\eqref{eq:K}。则下列结论成立。
\begin{enumerate}
\item[(i)] $\mathcal K_N$ 是 $\mathcal H_N$ 上的自伴希尔伯特--施密特算子。
\item[(ii)] 以 \eqref{eq:kW} 中的 $\kappa_W$ 与 $\kappa_{R_N}$，
\[
 \mathcal K_N=-W_N+R_N,\qquad W_N\ge0,\qquad W_N\ \text{是迹类的},
\]
且
\[
 \|R_N\|_{\HS}\ll(\log\log(3N))^2 .
\]
\item[(iii)] $\|W_N\|=\Bigl(\dfrac{\pi}{e^3}+o(1)\Bigr)\log N$。
\item[(iv)] $\mathcal K_N$ 的最小特征值满足
\[
 \lambda_{\min}(\mathcal K_N)=-\|W_N\|+O\bigl((\log\log N)^2\bigr)
 =-\Bigl(\frac{\pi}{e^3}+o(1)\Bigr)\log N,
\]
且 $\|\mathcal K_N\|=\bigl(\pi/e^3+o(1)\bigr)\log N$。
\item[(v)] $\|\mathcal K_N\|_{\HS}\ll(\log N)^{3/2}$。
\end{enumerate}
所有隐含常数是绝对的。
\end{theorem}

\begin{corollary}\label{cor:HS}
有 $\|\mathcal K_N\|_{\HS}\ll(\log N)^{3/2}$；特别地，\eqref{eq:split} 中模型
分量满足的希尔伯特--施密特界比 \cite[命题 6.5]{LL} 中使用的界
$\|\mathcal K_N\|_{\HS}\ll(\log N)^2$ 小一个因子 $(\log N)^{1/2}$。
\end{corollary}

\begin{remark}\label{rem:eigenvector}
\eqref{eq:w} 中的权重 $\varpi(x)=x-L_N+h_N$ 在 $[L_N,\infty)$ 上递增且无界，
故它在该区间上没有最大值。真正进入估计的是乘积 $\varpi(x)e^{-2\alpha_Nx}$；
由于
\[
 \bigl(\varpi(x)e^{-2\alpha_Nx}\bigr)'
 =e^{-2\alpha_Nx}\bigl(1-2\alpha_N\varpi(x)\bigr),
\]
该乘积在
\begin{equation}\label{eq:peakpoint}
 x^*=L_N+\frac1{2\alpha_N}-h_N
\end{equation}
处取得最大值，此时 $\varpi(x^*)=1/(2\alpha_N)$ 且
\begin{equation}\label{eq:peakvalue}
 \varpi(x^*)e^{-2\alpha_Nx^*}
 =\frac{e^{-2\alpha_NL_N-1+2\alpha_Nh_N}}{2\alpha_N}.
\end{equation}
引理~\ref{lem:twosided} 下界中的试验函数支撑在 $x^*$ 附近、宽度为
$\ll\alpha_N^{-1/2}$ 的窗口中；这就是“$W_N$ 的顶特征向量在 $x^*$ 附近集中”
的准确含义。
\end{remark}

\begin{remark}\label{rem:numerical}
对下文核 \eqref{eq:Gram} 作离散化，把
$\bigl|X_1(\alpha_N+it)/(1+\alpha_N+it)\bigr|^2$ 换成 $(1+t^2)^{-1}$，当
$\alpha_N=0.1$、$0.05$、$0.025$、$0.0125$ 时，最大特征值分别为
$\pi/(e^3\alpha_N)$ 的 $0.730$、$0.850$、$0.923$、$0.972$ 倍；该比值单调
趋于 $1$，与定理~\ref{thm:main}(iii) 相符。
\end{remark}

\subsection*{1.3.\ 方法}

证明分四步。

\begin{enumerate}
\item[\rm(A)] \emph{欧拉乘积。} 振幅 $\mathscr A_N(w)$ 改写为
$H_N(w)\zeta(1+w)\zeta(1-w)$，其中 $H_N=\mathfrak S(N)P_N/\zeta(1+w)$ 有显式
欧拉乘积且在 $\operatorname{Re}w\ge0$ 上解析；我们估计 $H_N$ 在原点处的值以及
前两阶导数（第~\ref{sec:amplitude} 节）。
\item[\rm(B)] \emph{极点分解。} $\zeta(1+w)\zeta(1-w)$ 在 $w=0$ 处的二阶极点
产生主部 $-w^{-2}-h_Nw^{-1}$；拉普拉斯恒等式把这一部分转化为负半定算子
$-W_N$，其核 \eqref{eq:kW} 是显式核 $\Phi_x$ 的积分（第~\ref{sec:laurent} 节）。
\item[\rm(C)] \emph{余项。} 余项核在原点附近与远离原点处分别估计；边界变量上的
二重积分给出 $\|R_N\|_{\HS}\ll(\log\log(3N))^2$（第~\ref{sec:remainder} 节）。
\item[\rm(D)] \emph{极点谱。} 通过两个变量 $x,y$ 下的核分析 $W_N$；与局部卷积核
的双边比较把 $\|W_N\|$ 确定到相对 $o(1)$，外尔不等式把结果转移到
$\mathcal K_N$（第~\ref{sec:polar}--\ref{sec:main} 节）。
\end{enumerate}

\subsection*{1.4.\ 约定}

记 $f\ll g$ 与 $f=O(g)$ 表示 $|f|\le Cg$，$C$ 为绝对常数；下标表示常数依赖。
除非另行说明，所有常数都是绝对的。我们不使用 $\mathrm{GRH}$ 或任何其他未证
假设；模型算子 $\mathcal K_N$ 是解析对象。

\section{预备}\label{sec:prelim}

\subsection*{2.1.\ 哈代空间}

我们回顾 \cite[\S5]{LL} 中 $\mathcal H_N$ 的描述。函数 $h$ 属于 $H^2(S_N)$，
如果它在 $S_N$ 中全纯且
\[
 \|h\|^2_{H^2(S_N)}
 =\sup_{0<\eps<(\beta-\alpha_N)/2}
 \int_{\mathbb R}\bigl(|h(\alpha_N+\eps+it)|^2+|h(\beta-\eps+it)|^2\bigr)\dd t<\infty .
\]
边界值在 $L^2$ 中存在，且 \eqref{eq:inner} 成立；两条竖线各带通常的勒贝格
测度。

\subsection*{2.2.\ 模型核}

核 \eqref{eq:kernel} 在条带内关于 $s$ 全纯，因为 $\mathscr A_N(w)$ 在
$\operatorname{Re}w>0$ 上全纯，而边界点满足
$\operatorname{Re}(s+\bar z)\ge2\alpha_N>0$。

\begin{lemma}\label{lem:selfadjoint}
对 $s,z\in\Gamma_N$ 有
\[
 \kappa_N(s,z)=\overline{\kappa_N(z,s)} .
\]
因此，一旦知道 $\mathcal K_N$ 是希尔伯特--施密特算子，它就是自伴的。
\end{lemma}

\begin{proof}
函数 $X_1$ 满足 $X_1(\bar z)=\overline{X_1(z)}$，而 $P_N$ 与 $\mathfrak S(N)$
的系数为实数，故对 $\operatorname{Re}w>0$ 有
$\mathscr A_N(\bar w)=\overline{\mathscr A_N(w)}$。于是
\[
 \begin{aligned}
 \overline{\kappa_N(z,s)}
 &=\frac{X_1(\bar z)}{1+\bar z}\cdot\frac{X_1(s)}{1+s}\,
 e^{-L_N(\bar z+s)}\overline{\mathscr A_N(z+\bar s)}\\
 &=\frac{X_1(s)X_1(\bar z)}{(1+s)(1+\bar z)}\,
 e^{-L_N(s+\bar z)}\mathscr A_N(s+\bar z)
 =\kappa_N(s,z).
 \end{aligned}
\]
\end{proof}

\begin{lemma}\label{lem:a}
设 $0<\rho\le\beta$，$t\in\mathbb R$。则
\begin{equation}\label{eq:abound}
 |X_1(\rho+it)|\asymp(1+|t|)^{-\rho},\qquad
 \int_{\mathbb R}\frac{|X_1(\rho+it)|^2}{|1+\rho+it|^2}\dd t\ll1,
\end{equation}
且对 $x\ge L_N$，
\begin{equation}\label{eq:phibound}
 \int_{\Gamma_N}\frac{|X_1(s)|^2}{|1+s|^2}\,e^{-2x\Re s}\,|\mathrm ds|
 \ll e^{-2\alpha_Nx}.
\end{equation}
\end{lemma}

\begin{proof}
斯特林公式给出 \eqref{eq:abound}；积分收敛是因为
$|X_1(\rho+it)|^2/|1+\rho+it|^2\asymp(1+|t|)^{-2\rho-2}$。把边界积分拆到两条
线上并用 \eqref{eq:abound} 即得 \eqref{eq:phibound}，因为 $\beta>\alpha_N$ 且
$x\ge L_N>0$。
\end{proof}

\subsection*{2.3.\ 延拓}

算子 \eqref{eq:K} 首先定义在 $\mathcal H_N$ 的稠密子空间上，该子空间由
傅里叶--拉普拉斯密度具紧支集的函数组成；在这个子空间上，下文中所有积分绝对
收敛。第~\ref{sec:remainder} 节的估计表明核在 $\Gamma_N\times\Gamma_N$ 上平方
可积，故 $\mathcal K_N$ 延拓为 $\mathcal H_N$ 上的希尔伯特--施密特算子；恒等式
$\mathcal K_N=-W_N+R_N$ 于是由连续性在 $\mathcal H_N$ 上成立。

\section{模型振幅}\label{sec:amplitude}

\subsection*{3.1.\ 一个欧拉乘积}

\begin{lemma}\label{lem:euler}
对 $\operatorname{Re}w>1$ 有
\begin{equation}\label{eq:HN}
 \begin{aligned}
 H_N(w):={}&\frac{\mathfrak S(N)P_N(w)}{\zeta(1+w)}\\
 ={}&2(1-2^{-1-w})\prod_{\substack{p\mid N\\p>2}}\frac{p}{p-1}\bigl(1-p^{-1-w}\bigr)\\
 &\times\prod_{\substack{p>2\\p\nmid N}}\Bigl(1-\frac1{(p-1)^2}\Bigr)
 \Bigl(1+\frac{p^{-w}}{p-2}\Bigr)\bigl(1-p^{-1-w}\bigr).
 \end{aligned}
\end{equation}
右端在 $\operatorname{Re}w\ge0$ 上局部一致收敛，并定义 $H_N$ 到该半平面的
解析延拓。
\end{lemma}

\begin{proof}
对 $\operatorname{Re}w>1$，
\[
 \frac1{\zeta(1+w)}=\prod_p(1-p^{-1-w})=(1-2^{-1-w})\prod_{p>2}(1-p^{-1-w}).
\]
代入 \eqref{eq:X1} 并把 $\prod_{p>2}$ 按 $p\mid N$ 与 $p\nmid N$ 拆开，
$p\mid N$、$p>2$ 处的因子为
\[
 \frac{p-1}{p-2}\Bigl(1-\frac1{(p-1)^2}\Bigr)\bigl(1-p^{-1-w}\bigr)
 =\frac{p}{p-1}\bigl(1-p^{-1-w}\bigr),
\]
而 $p\nmid N$、$p>2$ 处的因子为
\[
 \Bigl(1-\frac1{(p-1)^2}\Bigr)\Bigl(1+\frac{p^{-w}}{p-2}\Bigr)\bigl(1-p^{-1-w}\bigr).
\]
这就给出 \eqref{eq:HN}。对 $\operatorname{Re}w\ge0$ 与 $p\nmid N$，最后一个
因子等于 $1+O(p^{-2})$，关于 $w$ 一致，故乘积在该处局部一致收敛并定义解析函数；
它在 $\operatorname{Re}w>1$ 上与 $H_N$ 一致，故在整个半平面上一致。
\end{proof}

\subsection*{3.2.\ 原点处的值与一阶导数}

\begin{lemma}\label{lem:value}
$H_N$ 在 $\operatorname{Re}w\ge0$ 上解析，且
\begin{equation}\label{eq:HN0}
 H_N(0)=1,\qquad
 H_N'(0)=h_N:=\log2+\sum_{\substack{p\mid N\\p>2}}\frac{\log p}{p-1}.
\end{equation}
此外
\begin{equation}\label{eq:hNbound}
 0<h_N\ll\log\log(3N).
\end{equation}
\end{lemma}

\begin{proof}
在 $w=0$ 处，$2(1-2^{-1})=1$，$p\mid N$ 的有限乘积等于
$\prod p/(p-1)(1-1/p)=1$，而 $p\nmid N$ 的因子等于
\[
 \Bigl(1-\frac1{(p-1)^2}\Bigr)\Bigl(1+\frac1{p-2}\Bigr)\Bigl(1-\frac1p\Bigr)
 =\frac{p(p-2)}{(p-1)^2}\cdot\frac{p-1}{p-2}\cdot\frac{p-1}{p}=1 ;
\]
故 $H_N(0)=1$。对导数，我们在 $w=0$ 处对 \eqref{eq:HN} 取对数导数。第一个
因子贡献 $\log2\cdot2^{-1}/(1-2^{-1})=\log2$；$p\mid N$ 的因子贡献
$(\log p)p^{-1}/(1-p^{-1})=\log p/(p-1)$；$p\nmid N$ 的因子贡献
\[
 \begin{aligned}
 &-\frac{(\log p)p^{-w}}{p-2+p^{-w}}\Bigl|_{w=0}
 +\frac{(\log p)p^{-1-w}}{1-p^{-1-w}}\Bigl|_{w=0}\\
 &\qquad=-\frac{\log p}{p-1}+\frac{\log p}{p-1}=0 .
 \end{aligned}
\]
这就证明了 \eqref{eq:HN0}。对 \eqref{eq:hNbound}，我们用
$\prod_{p\mid N}p\le N$ 与标准估计
\[
 \sum_{p\mid N}\frac{\log p}{p}\ll\log\log(3N);
\]
$\log2$ 项无害。
\end{proof}

\subsection*{3.3.\ 导数界}

\begin{lemma}\label{lem:der}
对 $r=0,1,2$ 与 $\operatorname{Re}w\ge0$，
\begin{equation}\label{eq:derbound}
 |H_N^{(r)}(w)|\ll_r(\log\log(3N))^{r+2}.
\end{equation}
\end{lemma}

\begin{proof}
写 $H_N=\prod_pf_p$，局部因子为
\[
 f_2(w)=2(1-2^{-1-w}),\qquad
 f_p(w)=\frac{p}{p-1}\bigl(1-p^{-1-w}\bigr)\quad(p\mid N,\ p>2),
\]
\[
 f_p(w)=\Bigl(1-\frac1{(p-1)^2}\Bigr)\Bigl(1+\frac{p^{-w}}{p-2}\Bigr)
 \bigl(1-p^{-1-w}\bigr)\quad(p\nmid N,\ p>2).
\]
对 $\operatorname{Re}w\ge0$ 置 $z=p^{-w}$，则 $|z|\le1$。

\emph{局部估计。} 对 $p\nmid N$，恒等式
\[
 \Bigl(1+\frac{z}{p-2}\Bigr)\Bigl(1-\frac zp\Bigr)
 =1+\frac{2z-z^2}{p(p-2)}
\]
给出
\[
 f_p-1=\frac{2z-z^2}{p(p-2)}\Bigl(1-\frac1{(p-1)^2}\Bigr)-\frac1{(p-1)^2},
\]
故 $|f_p-1|\ll p^{-2}$；又因 $z'=-(\log p)z$，求导得
\[
 |f_p^{(r)}(w)|\ll_r(\log p)^rp^{-2}\qquad(r=1,2).
\]
对 $p\mid N$、$p>2$，有 $f_p=1+(1-z)/(p-1)$，故 $|f_p-1|\le2/(p-1)$，且
$r=1,2$ 时 $|f_p^{(r)}(w)|\ll_r(\log p)^r/(p-1)$。$p=2$ 的贡献及其导数为
$O(1)$。

\emph{求和。} 由于 $\prod_{p\mid N}p\le N$，
\begin{equation}\label{eq:sums}
 \begin{gathered}
 \sum_{p\mid N}\frac1{p-1}\ll\log\log\log(3N),\qquad
 \sum_{p\mid N}\frac{\log p}{p-1}\ll\log\log(3N),\\
 \sum_{p\mid N}\frac{(\log p)^2}{p-1}\ll(\log\log(3N))^2.
 \end{gathered}
\end{equation}
且 $\sum_{p}(\log p)^jp^{-2}\ll_j1$。因此
\[
 |H_N|\le\prod_p(1+|f_p-1|)\le\exp\Bigl(2\sum_{p\mid N}\frac1{p-1}+O(1)\Bigr)\ll(\log\log(3N))^2 .
\]
对导数，我们直接对乘积求导而不做除法：
\[
 H_N'=\sum_pf_p'\prod_{q\ne p}f_q,\qquad
 H_N''=\sum_pf_p''\prod_{q\ne p}f_q+\sum_{p\ne q}f_p'f_q'\prod_{r\ne p,q}f_r .
\]
对每个 $p$，$\prod_{q\ne p}|f_q|\le\prod_{q\ne p}(1+|f_q-1|)\ll(\log\log(3N))^2$，且由
\eqref{eq:sums} 有 $\sum_p|f_p'|\ll\log\log(3N)$、$\sum_p|f_p''|\ll(\log\log(3N))^2$。因此
\[
 \begin{aligned}
 |H_N'|&\ll(\log\log(3N))^2\sum_p|f_p'|\ll(\log\log(3N))^3,\\
 |H_N''|&\ll(\log\log(3N))^2\Bigl(\sum_p|f_p''|+\bigl(\sum_p|f_p'|\bigr)^2\Bigr)
 \ll(\log\log(3N))^4 .
 \end{aligned}
\]
这就证明了 \eqref{eq:derbound}。
\end{proof}

\begin{remark}\label{rem:zero}
关于 $p\nmid N$ 的乘积可能在虚轴上的孤立点为零（例如 $p=3$、
$w=(2k+1)\pi i/\log3$ 处），但上述估计通过逐项对乘积求导得到，不需要对数
导数。
\end{remark}

\begin{remark}\label{rem:sharp}
\eqref{eq:derbound} 中的指数是最优的：取 $N=\prod_{p\le y}p$，则
$\log\log(3N)\asymp\log y$；由 $\{t\log p\bmod2\pi\}_{p\le y}$ 在环面上的稠密性，存在
$t$ 使所有因子 $|1-p^{-1-it}|$ 都接近 $1+p^{-1}$；在这些点上
$|H_N(it)|\asymp\prod_{p\le y}(p+1)/(p-1)\asymp(\log\log(3N))^2$，同理
$|H_N'|\asymp(\log\log(3N))^3$、$|H_N''|\asymp(\log\log(3N))^4$。在原点处数值要小得多：
$H_N(0)=1$、$H_N'(0)=h_N\ll\log\log(3N)$，且
\[
 H_N''(0)=h_N^2+\sum_{\substack{p\mid N\\p>2}}\frac{(\log p)^2(p-2)}{(p-1)^2}+O(1)\ll(\log\log(3N))^2 .
\]
\end{remark}

\section{原点处的洛朗展开}\label{sec:laurent}

\subsection*{4.1.\ $\zeta$ 乘积}

\begin{lemma}\label{lem:zeta}
在 $w=0$ 的邻域内，
\begin{equation}\label{eq:zetaprod}
 \zeta(1+w)\zeta(1-w)=-\frac1{w^2}+\gamma^2+2\gamma_1+O(w^2),
\end{equation}
其中 $\gamma$ 是欧拉常数，$\gamma_1$ 是第一个斯蒂尔切斯常数。
\end{lemma}

\begin{proof}
函数 $w\mapsto\zeta(1+w)\zeta(1-w)$ 是偶函数，故其关于 $w=0$ 的洛朗展开只含
$w$ 的偶次幂。洛朗展开
\[
 \zeta(1+w)=\frac1w+\gamma-\gamma_1w+O(w^2),\qquad
 \zeta(1-w)=-\frac1w+\gamma+\gamma_1w+O(w^2)
\]
相乘即得 \eqref{eq:zetaprod}，其中奇次项为零。
\end{proof}

\subsection*{4.2.\ 主部}

\begin{proposition}\label{prop:laurent}
存在绝对常数 $\delta>0$，使得对 $|w|\le2\delta$、$\operatorname{Re}w\ge0$，
\begin{equation}\label{eq:laurent}
 \mathscr A_N(w)=-\frac1{w^2}-\frac{h_N}{w}+R_N^{(0)}(w),
\end{equation}
其中 $R_N^{(0)}$ 在该半圆盘上全纯，且
\begin{equation}\label{eq:R0}
 R_N^{(0)}(w)=c_0+O\bigl((\log\log(3N))^4\bigr),
 \qquad c_0=\gamma^2+2\gamma_1,
\end{equation}
特别地
\begin{equation}\label{eq:R0near}
 |R_N^{(0)}(w)|\ll(\log\log(3N))^4 .
\end{equation}
\end{proposition}

\begin{proof}
由 \eqref{eq:HN} 与引理~\ref{lem:zeta}，
\[
 \mathscr A_N(w)=H_N(w)\zeta(1+w)\zeta(1-w).
\]
由引理~\ref{lem:der} 取 $r=2$，并对从 $0$ 到 $w$ 的线段（它落在
$\operatorname{Re}w\ge0$ 内）用泰勒公式，
\[
 H_N(w)=1+h_Nw+O(w^2(\log\log(3N))^4),\qquad |w|\le2\delta,\quad\operatorname{Re}w\ge0 .
\]
乘以 \eqref{eq:zetaprod} 并减去两个极点项，即得 \eqref{eq:laurent}，其中
\[
 R_N^{(0)}(w)=c_0+O((\log\log(3N))^4)+O(w)+O(w^2(\log\log(3N))^4),
\]
它全纯；这里的第一个误差来自 $O(w^2(\log\log(3N))^4)\cdot(-w^{-2})$。这就给出
\eqref{eq:R0}。
\end{proof}

\subsection*{4.3.\ 拉普拉斯表示}

\begin{lemma}\label{lem:laplace}
对 $\operatorname{Re}w>0$，
\begin{equation}\label{eq:laplace1}
 \frac{e^{-L_Nw}}{w^2}=\int_{L_N}^{\infty}(x-L_N)e^{-xw}\dd x,
 \qquad
 \frac{e^{-L_Nw}}{w}=\int_{L_N}^{\infty}e^{-xw}\dd x .
\end{equation}
\end{lemma}

\begin{proof}
在第一个积分中作代换 $u=x-L_N$ 得
$e^{-L_Nw}\int_0^\infty ue^{-uw}\dd u=e^{-L_Nw}/w^2$；第二个是初等的。
\end{proof}

\subsection*{4.4.\ 极点算子}

\begin{proposition}\label{prop:polar}
\eqref{eq:kW} 中的算子 $W_N$ 是有定义的、正且迹类的，并且
\begin{equation}\label{eq:tracepole}
 \tr W_N=\int_{L_N}^{\infty}\varpi(x)\int_{\Gamma_N}\frac{|X_1(s)|^2}{|1+s|^2}
 e^{-2x\Re s}\,|\mathrm ds|\,\mathrm dx<\infty .
\end{equation}
此外 $W_N=-\mathcal K_N^{pole}$，其中 $\mathcal K_N^{pole}$ 的核是 $\kappa_N$ 在 $w=0$ 处的
主部。
\end{proposition}

\begin{proof}
对 $s,z\in\Gamma_N$ 置 $w=s+\bar z$，则 $\operatorname{Re}w\ge2\alpha_N>0$。
由 \eqref{eq:laurent} 与引理~\ref{lem:laplace}，$\kappa_N$ 的主部等于
\[
 \frac{X_1(s)X_1(\bar z)}{(1+s)(1+\bar z)}e^{-L_Nw}
 \Bigl(-\frac1{w^2}-\frac{h_N}{w}\Bigr)
 =-\int_{L_N}^{\infty}\varpi(x)\,\Phi_x(s,z)\dd x=-\kappa_W(s,z).
\]
故 $\mathcal K_N^{pole}$ 的核为 $-\kappa_W$，即 $W_N=-\mathcal K_N^{pole}$。对
$h\in\mathcal H_N$，$W_N$ 的二次型为
\[
 \langle W_Nh,h\rangle
 =\int_{L_N}^{\infty}\varpi(x)\Bigl|\int_{\Gamma_N}
 \frac{h(s)X_1(\bar s)}{1+\bar s}e^{-\bar s x}\,|\mathrm ds|\Bigr|^2\dd x\ge0,
\]
故 $W_N\ge0$。迹公式 \eqref{eq:tracepole} 由同一核沿对角线积分得到；由
\eqref{eq:phibound} 它有限，因为
$\int_{L_N}^{\infty}\varpi(x)e^{-2\alpha_Nx}\dd x<\infty$。
\end{proof}

\section{余项}\label{sec:remainder}

\subsection*{5.1.\ 余项核}

由 \eqref{eq:laurent}，差 $\mathcal K_N+W_N$ 的核为
\begin{equation}\label{eq:Rkernel}
 \kappa_{R_N}(s,z)=\frac{X_1(s)X_1(\bar z)}{(1+s)(1+\bar z)}\,
 e^{-L_N(s+\bar z)}R_N^{(0)}(s+\bar z),
\end{equation}
其中 $R_N^{(0)}$ 是命题~\ref{prop:laurent} 的全纯函数。

\subsection*{5.2.\ 远离原点}

\begin{lemma}\label{lem:far}
设 $\eps>0$ 且 $|w|\ge2\delta$、
$0\le\operatorname{Re}w\le2\beta=\frac14$。则
\begin{equation}\label{eq:far}
 |R_N^{(0)}(w)|\ll_{\delta,\eps}
 (\log\log(3N))^2(2+|\Im w|)^{1/8+\eps}\log(2+|\Im w|).
\end{equation}
\end{lemma}

\begin{proof}
对 $|w|\ge2\delta$ 有
$R_N^{(0)}(w)=\mathscr A_N(w)+w^{-2}+h_Nw^{-1}$，后两项为
$O_\delta(1+h_N)=O_\delta(\log\log(3N))$。只需估计
$\mathscr A_N=H_N\zeta(1+w)\zeta(1-w)$。

由引理~\ref{lem:der} 取 $r=0$，$|H_N(w)|\ll(\log\log(3N))^2$。对因子
$\zeta(1+w)$ 分两种情形。
若 $\Re w\ge\delta/2$，则 $\operatorname{Re}(1+w)\ge1+\delta/2>1$，且
$|\zeta(1+w)|\le\zeta(1+\delta/2)\ll_\delta1$。若 $\Re w<\delta/2$，则
$|w|\ge2\delta$ 迫使 $|\Im w|\ge\sqrt3\delta/2$，而对半平面
$\operatorname{Re}s\ge1$ 的标准估计给出
$|\zeta(1+w)|\ll_\delta\log(2+|\Im w|)$。对因子 $\zeta(1-w)$，
$1-\Re w\in[\frac34,1)$ 的凸性界给出
\[
 |\zeta(1-w)|\ll_\eps(2+|\Im w|)^{\Re w/2+\eps}\le(2+|\Im w|)^{1/8+\eps}.
\]
这就证明了 \eqref{eq:far}。
\end{proof}

\subsection*{5.3.\ 原点附近}

\begin{lemma}\label{lem:near}
对 $|w|\le2\delta$、$\operatorname{Re}w\ge0$，有
\begin{equation}\label{eq:near}
 |R_N^{(0)}(w)|\ll_\delta(\log\log(3N))^2 .
\end{equation}
\end{lemma}

\begin{proof}
记 $L=\log\log(3N)$。

\emph{圆弧。} 设 $|w|=2\delta$、$\operatorname{Re}w\ge0$。由引理~\ref{lem:far}，
\[
 |R_N^{(0)}(w)|\ll_{\delta,\eps}L^2(2+|\Im w|)^{1/8+\eps}\log(2+|\Im w|)
 \ll_\delta L^2 .
\]

\emph{直径。} 设 $f_p$ 为引理~\ref{lem:der} 的局部因子。在半圆盘
$\{|w|\le2\delta,\ \operatorname{Re}w\ge0\}$ 上每个 $f_p$ 非零——因子的零点都在
其外（$p\mid N$ 时零点满足 $|p^{-w}|=p$；$\operatorname{Re}w\ge0$ 上的零点只可能
来自 $p=3$，位于 $w=(2k+1)\pi i/\log3$，$|w|>\pi/\log3$）——故存在
$\log H_N=\sum_p\log f_p$ 的全纯分支 $g$，且 $g(0)=0$。于是 $g'(0)=h_N$，且
\[
 |g''(w)|\ll\sum_p\Bigl|\frac{f_p''}{f_p}-\Bigl(\frac{f_p'}{f_p}\Bigr)^2\Bigr|
 \ll L^2,
\]
因为 $p\mid N$ 时 $|\partial_w^2\log f_p(w)|\ll(\log p)^2/p$，$p\nmid N$ 时
$\ll(\log p)^2/p^2$，而 $p=2$ 贡献 $O(1)$。沿从 $0$ 到 $it$ 的线段用泰勒公式
（余项写成积分形式，并利用 $g''$ 在 $\operatorname{Re}w\ge0$ 上的连续性），得
\[
 g(it)=ith_N+O(t^2L^2),\qquad\text{从而}\qquad
 H_N(it)=1+ith_N+O(t^2L^2).
\]
由引理~\ref{lem:zeta}，对 $Z(w)=\zeta(1+w)\zeta(1-w)$ 有
$Z(it)=t^{-2}+c_0+O(t^2)$；又由引理~\ref{lem:der} 取 $r=0$ 得
$|H_N(it)|\ll L^2$。由 \eqref{eq:laurent}，
$R_N^{(0)}=H_NZ+w^{-2}+h_Nw^{-1}$，故对 $0<|t|\le2\delta$ 有
\[
 R_N^{(0)}(it)=\frac{H_N(it)-1-ith_N}{t^2}+c_0H_N(it)+O(L^2t^2)=O(L^2),
\]
而 $t=0$ 处由连续性得同样界。

\emph{结论。} $R_N^{(0)}$ 在半圆盘
$\{|w|<2\delta,\ \operatorname{Re}w>0\}$ 上全纯、在其闭包上连续；圆弧与直径上
都界以 $C_\delta L^2$，故由最大模原理得 \eqref{eq:near}。
\end{proof}

\subsection*{5.4.\ 希尔伯特--施密特估计}

\begin{proposition}\label{prop:HS}
对每个 $0<\eps<3/8$，
\begin{equation}\label{eq:HS}
 \|R_N\|_{\HS}^2\ll_\eps(\log\log(3N))^4
 \iint_{\mathbb R^2}
 \frac{(2+|t-u|)^{1/4+2\eps}\log^2(2+|t-u|)}
      {(1+t^2)(1+u^2)}\,\mathrm dt\,\mathrm du ,
\end{equation}
从而
\begin{equation}\label{eq:HSbound}
 \|R_N\|_{\HS}\ll(\log\log(3N))^2 .
\end{equation}
\end{proposition}

\begin{proof}
由 \eqref{eq:Rkernel}，
\[
 \|R_N\|_{\HS}^2
 =\int_{\Gamma_N}\int_{\Gamma_N}
 \frac{|X_1(s)|^2|X_1(z)|^2}{|1+s|^2|1+\bar z|^2}\,
 e^{-2L_N\operatorname{Re}(s+\bar z)}
 |R_N^{(0)}(s+\bar z)|^2\,|\mathrm ds|\,|\mathrm dz| .
\]
由于 $\operatorname{Re}(s+\bar z)\ge2\alpha_N$ 且 $L_N>0$，指数至多为 $1$。
写 $s=\rho+it$、$z=\rho'+iu$，$\rho,\rho'\in\{\alpha_N,\beta\}$。近原点区域
$|s+\bar z|\le2\delta$ 由 \eqref{eq:near} 与 \eqref{eq:abound} 贡献
$O_\delta((\log\log(3N))^4)$。在其余区域，\eqref{eq:far} 适用（$\Im w=t-u$），
而 \eqref{eq:abound} 给出因子 $(1+t^2)^{-1}(1+u^2)^{-1}$。这给出
\eqref{eq:HS}。

还需验证二重积分收敛。取 $0<\eps<3/8$，置 $\theta=\frac14+2\eps<1$。由于
$(2+|t-u|)^\theta\ll_\theta(2+|t|)^\theta(2+|u|)^\theta$，有
\[
 \iint_{\mathbb R^2}\frac{(2+|t-u|)^\theta\log^2(2+|t-u|)}
      {(1+t^2)(1+u^2)}\dd t\dd u
 \ll_\theta\Bigl(\int_{\mathbb R}\frac{(2+|t|)^\theta\log^2(2+|t|)}{1+t^2}\dd t\Bigr)^2<\infty,
\]
收敛性由 $\theta<1$ 得出。这就证明了 \eqref{eq:HSbound}。
\end{proof}

\begin{proof}[定理~\ref{thm:main}(i)--(ii) 证明的完成]
恒等式 $\mathcal K_N=-W_N+R_N$ 是由 \eqref{eq:laurent} 结合 \eqref{eq:kW}
与 \eqref{eq:Rkernel} 得到的核恒等式。由于 $W_N$ 由命题~\ref{prop:polar}
是迹类的、$R_N$ 由命题~\ref{prop:HS} 是希尔伯特--施密特的，$\mathcal K_N$ 是
希尔伯特--施密特算子；其自伴性由引理~\ref{lem:selfadjoint} 给出。这就证明了
(i) 与 (ii)。
\end{proof}

\section{极点算子}\label{sec:polar}

\subsection*{6.1.\ 两变量下的核}

对 $x,y\ge L_N$ 置
\begin{equation}\label{eq:Gram}
 \mathscr G(x,y)=\int_{\Gamma_N}\frac{X_1(s)X_1(\bar s)}{|1+s|^2}
 e^{-xs-\bar sy}\,|\mathrm ds| .
\end{equation}
在两条线上代入 $s=\rho+it$ 得
\begin{equation}\label{eq:g}
 \mathscr G(x,y)=\sum_{\rho\in\{\alpha_N,\beta\}}e^{-\rho(x+y)}\,g_\rho(x-y),
\end{equation}
其中
\begin{equation}\label{eq:grho}
 g_\rho(v)=\int_{\mathbb R}\frac{|X_1(\rho+it)|^2}{|1+\rho+it|^2}
 e^{-itv}\dd t,\qquad \rho\in\{\alpha_N,\beta\}.
\end{equation}

\begin{lemma}\label{lem:gram}
每个 $g_\rho$ 是实值、偶的、正定的，且
\begin{equation}\label{eq:gmass}
 \int_{\mathbb R}g_\rho(v)\dd v=2\pi\,\frac{|X_1(\rho)|^2}{|1+\rho|^2} .
\end{equation}
此外，对 $h\in\mathcal H_N$，
\begin{equation}\label{eq:qform}
 \langle W_Nh,h\rangle
 =\int_{L_N}^{\infty}\varpi(x)\Bigl|\int_{\Gamma_N}
 \frac{h(s)X_1(\bar s)}{1+\bar s}e^{-\bar s x}\,|\mathrm ds|\Bigr|^2\dd x,
\end{equation}
且 $\|W_N\|$ 等于 $L^2([L_N,\infty))$ 上以
\begin{equation}\label{eq:Wkernel}
 \mathscr W_N(x,y)=\sqrt{\varpi(x)\varpi(y)}\,\mathscr G(x,y)
\end{equation}
为核的积分算子的范数。
\end{lemma}

\begin{proof}
在 \eqref{eq:Gram} 中代入 $s=\rho+it$ 即得 \eqref{eq:g}。对固定的 $\rho$，
函数 $t\mapsto|X_1(\rho+it)|^2/|1+\rho+it|^2$ 是偶的非负函数，故 $g_\rho$ 是
实值、偶的、正定的；\eqref{eq:gmass} 是 $v=0$ 处的傅里叶反演公式。
\eqref{eq:qform} 是核 $\kappa_W$ 的二次型。

对最后一个断言，令 $V:L^2([L_N,\infty))\to\mathcal H_N$ 为
\[
 (Vf)(s)=\int_{L_N}^{\infty}\sqrt{\varpi(x)}\,f(x)\,
 \frac{X_1(s)}{1+s}e^{-sx}\dd x .
\]
其伴随为
\[
 (V^*h)(x)=\sqrt{\varpi(x)}\int_{\Gamma_N}h(s)\frac{X_1(\bar s)}{1+\bar s}
 e^{-\bar s x}\,|\mathrm ds| .
\]
于是
\[
 \langle VV^*h,h\rangle=\|V^*h\|^2
 =\int_{L_N}^{\infty}\varpi(x)\Bigl|\int_{\Gamma_N}h(s)\frac{X_1(\bar s)}{1+\bar s}
 e^{-\bar s x}\,|\mathrm ds|\Bigr|^2\dd x,
\]
即 \eqref{eq:qform}；故 $VV^*$ 就是算子 $W_N$。另一方面，$V^*V$ 的核是
\eqref{eq:Wkernel} 的共轭；共轭核对应的算子是原算子的伴随，故范数相同。因此
$\|W_N\|=\|VV^*\|=\|V^*V\|=\|\mathscr W_N\|$。
\end{proof}

\subsection*{6.2.\ 双边范数引理}

\begin{lemma}\label{lem:twosided}
设 $0<\alpha\le\alpha_0$，$L\ge0$，$\eta\ge0$，在 $[L,\infty)$ 上
$\varpi(x)=x-L+\eta$，$D(x)=\sqrt{\varpi(x)}e^{-\alpha x}$。设 $g\in L^1(\mathbb R)$ 是偶的、
正定的，且 $\int_{\mathbb R}|v||g(v)|\dd v<\infty$，$M$ 是 $L^2(\mathbb R)$ 上以
$g$ 为核的卷积算子，故 $\|M\|=\int g$。假设
$x\mapsto \varpi(x)e^{-2\alpha x}$ 在某个 $x^*>L$ 处取得最大值。则
\begin{equation}\label{eq:twosided}
 \|DMD\|=\bigl(1+O(\alpha^{1/2})\bigr)\|M\|\max_{x\ge L}\bigl(\varpi(x)e^{-2\alpha x}\bigr).
\end{equation}
\end{lemma}

\begin{proof}
\emph{上界。} $D$ 有界，故
$\|DMD\|\le\|D\|_\infty^2\|M\|=\|M\|\max \varpi(x)e^{-2\alpha x}$。

\emph{下界。} 置 $B=\alpha^{-1/2}$，令
$f=B^{-1/2}\mathbf 1_{[x^*-B/2,x^*+B/2]}$，故 $\|f\|_2=1$。记
$F(x)=\varpi(x)e^{-2\alpha x}$。由于 $F'(x^*)=0$ 且
$F''(x^*)/F(x^*)=-1/\varpi(x^*)^2$，泰勒公式给出
\[
 F(x)\ge F(x^*)\bigl(1-O(\alpha^2(x-x^*)^2)\bigr)\ge F(x^*)(1-O(\alpha)),
 \qquad |x-x^*|\le B/2 .
\]
在 $f$ 的支集上有 $D(x)^2=F(x)$。故
\[
 \langle DMDf,f\rangle
 =\iint g(x-y)D(x)f(x)D(y)f(y)\dd x\dd y
 \ge F(x^*)(1-O(\alpha))\langle Mf,f\rangle .
\]
对最后一个因子，
\[
 \begin{split}
 \langle Mf,f\rangle
 &=\frac1B\int_0^B2(B-v)g(v)\dd v
 =2\int_0^Bg(v)\dd v-\frac2B\int_0^Bvg(v)\dd v\\
 &=\|M\|+O\Bigl(\int_B^\infty|g|+\frac1B\int_0^\infty v|g|\Bigr).
 \end{split}
\]
由于 $g$ 可积且一阶矩有限，误差为 $O(B^{-1})=O(\alpha^{1/2})$。因此
$\langle DMDf,f\rangle\ge(1-O(\alpha^{1/2}))\|M\|F(x^*)$。
\end{proof}

\subsection*{6.3.\ $W_N$ 的范数}

\begin{proposition}\label{prop:norm}
有
\begin{equation}\label{eq:normW}
 \|W_N\|=\Bigl(\frac{\pi}{e^3}+o(1)\Bigr)\frac1{\alpha_N}
 =\Bigl(\frac{\pi}{e^3}+o(1)\Bigr)\log N .
\end{equation}
\end{proposition}

\begin{proof}
由 \eqref{eq:g} 与 \eqref{eq:Wkernel}，核 $\mathscr W_N$ 分成
$\alpha_N$ 部分与 $\beta$ 部分；两个核都是正的，故
\[
 \|W_N\|\ge\|\mathscr W_{\alpha_N}\|,\qquad
 \|W_N\|\le\|\mathscr W_{\alpha_N}\|+\|\mathscr W_\beta\| .
\]

\emph{$\alpha_N$ 线。} 这里
$\mathscr W_{\alpha_N}(x,y)=D(x)g_{\alpha_N}(x-y)D(y)$，其中
$D(x)=\sqrt{\varpi(x)}e^{-\alpha_Nx}$。由 \eqref{eq:grho} 与傅里叶反演公式，
\[
 \|M_{\alpha_N}\|=2\pi\sup_t\frac{|X_1(\alpha_N+it)|^2}{|1+\alpha_N+it|^2}
 =2\pi\frac{|X_1(\alpha_N)|^2}{(1+\alpha_N)^2}(1+O(\alpha_N))=2\pi(1+O(\alpha_N)),
\]
其中 $M_{\alpha_N}$ 是以 $g_{\alpha_N}$ 为核的卷积。右端的最大值在 $t=0$ 处
取得：在 $|t|\le2$ 上比值等于 $t=0$ 处的值乘以
$1+O(\alpha_N\log(2+|t|))$，而在 $|t|>2$ 上它 $\ll t^{-2}$，当 $\alpha_N$ 充分
小时比 $t=0$ 处的值小一个绝对因子。

乘积 $x\mapsto \varpi(x)e^{-2\alpha_Nx}$ 在 \eqref{eq:peakpoint} 的 $x^*$ 处取得
最大值，值为 \eqref{eq:peakvalue}，且对大 $N$ 有 $x^*>L_N$。取
$\alpha=\alpha_N$、$\eta=h_N$、$g=g_{\alpha_N}$ 应用引理~\ref{lem:twosided} 得
\[
 \|\mathscr W_{\alpha_N}\|
 =(1+O(\alpha_N^{1/2}))\,2\pi\,\frac{|X_1(\alpha_N)|^2}{(1+\alpha_N)^2}\,
 \frac{e^{-2\alpha_NL_N-1+2\alpha_Nh_N}}{2\alpha_N}.
\]

\emph{$\beta$ 线。} 类似地，
\[
 \|\mathscr W_\beta\|\le\|M_\beta\|\max_{x\ge L_N}\varpi(x)e^{-2\beta x}\ll N^{-1/4},
\]
它是 $o(\alpha_N^{-1})$。

\emph{结论。} 由于 $e^{-2\alpha_NL_N}=e^{-2}(1+O(\alpha_N))$、
$e^{2\alpha_Nh_N}=1+O(\alpha_Nh_N)=1+o(1)$ 且
$|X_1(\alpha_N)|^2=1+O(\alpha_N)$，我们得到
\[
 \|W_N\|=(1+o(1))\,\frac{\pi}{e^3}\cdot\frac1{\alpha_N},
\]
即 \eqref{eq:normW}。
\end{proof}

\subsection*{6.4.\ 迹与希尔伯特--施密特范数}

\begin{lemma}\label{lem:trace}
有
\begin{equation}\label{eq:trace}
 \tr W_N\ll\alpha_N^{-2},\qquad
 \|W_N\|_{\HS}\ll\alpha_N^{-3/2}.
\end{equation}
\end{lemma}

\begin{proof}
由 \eqref{eq:tracepole} 与 \eqref{eq:phibound}，
\[
 \tr W_N=\int_{L_N}^{\infty}\varpi(x)\int_{\Gamma_N}\frac{|X_1(s)|^2}{|1+s|^2}
 e^{-2x\Re s}|\mathrm ds|\,\mathrm dx\ll\int_{L_N}^{\infty}\varpi(x)e^{-2\alpha_Nx}\dd x
 \ll\alpha_N^{-2}.
\]
由于 $W_N\ge0$，
$\|W_N\|_{\HS}^2=\tr(W_N^2)\le\|W_N\|\tr W_N\ll\alpha_N^{-3}$。
\end{proof}

\section{最小特征值}\label{sec:main}

\subsection*{7.1.\ 外尔不等式}

\begin{lemma}\label{lem:weyl}
设 $A$ 是希尔伯特空间上自伴紧算子，$B$ 是自伴希尔伯特--施密特算子。则
\[
 \bigl|\lambda_{\min}(A+B)-\lambda_{\min}(A)\bigr|\le\|B\|_{\HS}.
\]
\end{lemma}

\begin{proof}
这是极小极大原理连同 $\|f\|=1$ 时的
$|\langle Bf,f\rangle|\le\|B\|_{\HS}$。
\end{proof}

\subsection*{7.2.\ 定理~\ref{thm:main} 的证明}

我们对 $A=-W_N$、$B=R_N$ 应用引理~\ref{lem:weyl}。由于 $W_N\ge0$ 是紧的
迹类算子，$\lambda_{\min}(-W_N)=-\|W_N\|$。由命题~\ref{prop:HS}，
$\|R_N\|_{\HS}\ll(\log\log(3N))^2$。故
\[
 \lambda_{\min}(\mathcal K_N)=-\|W_N\|+O\bigl((\log\log N)^2\bigr),
\]
而命题~\ref{prop:norm} 给出 (iv)。对算子范数，我们用
\[
 \bigl|\,\|\mathcal K_N\|-\|W_N\|\,\bigr|\le\|R_N\|\le\|R_N\|_{\HS},
\]
故 $\|\mathcal K_N\|=(\pi/e^3+o(1))\log N$。最后，
\[
 \|\mathcal K_N\|_{\HS}\le\|W_N\|_{\HS}+\|R_N\|_{\HS}
 \ll\alpha_N^{-3/2}+(\log\log(3N))^2\ll(\log N)^{3/2},
\]
即 (v)。证毕。

\subsection*{7.3.\ 注}

\begin{remark}\label{rem:half}
定理~\ref{thm:main}(v) 中的指数 $3/2$ 来自
$\|W_N\|_{\HS}^2\le\|W_N\|\tr W_N$，其中
$\|W_N\|\ll\alpha_N^{-1}$、$\tr W_N\ll\alpha_N^{-2}$；它未必最优。
\end{remark}

\begin{remark}\label{rem:noGRH}
证明只用到 $\mathscr A_N$ 在条带 $\operatorname{Re}w\ge2\alpha_N>0$ 上的解析
性质、$1/\zeta$ 的欧拉乘积以及斯特林公式。特别地，定理~\ref{thm:main} 是
无条件的；假设 $\mathrm{GRH}(N)$ 只通过分解 \eqref{eq:split} 与特征分量
$\mathcal P_N$ 进入 \cite{LL} 的哥德巴赫论证。
\end{remark}

\begin{remark}\label{rem:kernel-form}
\eqref{eq:Phi} 的核 $\Phi_x$ 是 \cite{LL} 的标准边界因子乘以指数
$e^{-x(s+\bar z)}$；引理~\ref{lem:laplace} 的拉普拉斯恒等式是唯一用到这个指数
的地方。特别地，算子 $W_N$ 由它的核 $\kappa_W$ 确定，所有估计都是核的估计。
\end{remark}

\begin{thebibliography}{9}

\bibitem[LL]{LL}
J. Li and J. Liu,
\emph{A conditional proof of the Goldbach conjecture},
arXiv:2606.05224 (2026).

\end{thebibliography}

\end{document}
